\title{Επικουρικές Υπηρεσίες Σχετικές με τη Συχνότητα}

\section{Εισαγωγή}
\label{kef-03-theoria}

Η συχνότητα ενός συστήματος ηλεκτρικής ενέργειας παραμένει σταθερή μόνο όσο η παραγωγή ισούται
σε κάθε στιγμή με τη ζήτηση και τις απώλειες. Οι ανισορροπίες που απομένουν μετά τον
προγραμματισμό της λειτουργίας και τις αγορές (Κεφάλαιο 1, ενότητα «\nameref{sec:stages}»),
όπως τα σφάλματα πρόβλεψης φορτίου και ανανεώσιμης παραγωγής, οι βηματικές μεταβολές των
προγραμμάτων και οι απώλειες μονάδων ηλεκτροπαραγωγής, καλύπτονται από τις \emph{επικουρικές
υπηρεσίες σχετικές με τη συχνότητα}. Ο Διαχειριστής Συστήματος Μεταφοράς (ΔΣΜ) δεσμεύει τις
αντίστοιχες εφεδρείες ενεργού ισχύος από τους παρόχους υπηρεσιών εξισορρόπησης και τις
ενεργοποιεί αυτόματα ή με εντολές κατανομής.

Τις κατηγορίες εφεδρειών και τις παραμέτρους ποιότητας συχνότητας κάθε συγχρονισμένης περιοχής
ορίζει ο Κανονισμός (ΕΕ) 2017/1485 για τη λειτουργία του συστήματος μεταφοράς (System
Operation Guideline, SO GL), ενώ τη δέσμευση και την ενεργοποίησή τους μέσω αγορών ρυθμίζει ο
Κανονισμός (ΕΕ) 2017/2195 για την εξισορρόπηση (Electricity Balancing Guideline, EB GL). Στην
Ελλάδα οι εφεδρείες δεσμεύονται και ενεργοποιούνται μέσω της Αγοράς Εξισορρόπησης του ΑΔΜΗΕ
(Κεφάλαιο 1, υποενότητα «\nameref{sec:greek-market}»).

Το κεφάλαιο εξετάζει αρχικά τις κατηγορίες εφεδρειών συχνότητας και στη συνέχεια τους
μηχανισμούς στους οποίους στηρίζονται: την αδράνεια των σύγχρονων μηχανών, την εφεδρεία
διατήρησης συχνότητας με τον στατισμό των ρυθμιστών στροφών και την εφεδρεία αποκατάστασης
συχνότητας, η οποία υλοποιείται από την Αυτόματη Ρύθμιση Παραγωγής (ΑΡΠ· Automatic Generation
Control, AGC), σε απομονωμένο και σε διασυνδεδεμένο σύστημα. Κλείνει με τις παραμέτρους
συχνότητας της Ηπειρωτικής Ευρώπης.

\section{Κατηγορίες εφεδρειών συχνότητας}
\label{sec:reserves}

Μετά από ανισορροπία μεταξύ παραγωγής και ζήτησης, όπως η απώλεια μονάδας ηλεκτροπαραγωγής, η
συχνότητα ελέγχεται από μια ακολουθία μηχανισμών με διαδοχικά βραδύτερη απόκριση. Αμέσως μετά
τη διαταραχή η ανισορροπία καλύπτεται από την κινητική ενέργεια των στρεφόμενων μαζών των
σύγχρονων μηχανών, χωρίς ενέργεια ελέγχου. Ακολουθούν οι εφεδρείες ενεργού ισχύος που ορίζει ο
SO GL:

\begin{itemize}
\item \textbf{Εφεδρεία διατήρησης συχνότητας} (ΕΔΣ· Frequency Containment Reserve, FCR).
      Αυτόματη και τοπική, ενεργοποιείται από τους ρυθμιστές στροφών των μονάδων από κοινού σε
      ολόκληρη τη συγχρονισμένη περιοχή. Συγκρατεί και σταθεροποιεί τη συχνότητα, με μόνιμη
      απόκλιση ως προς την ονομαστική τιμή.
\item \textbf{Εφεδρεία αποκατάστασης συχνότητας} (ΕΑΣ· Frequency Restoration Reserve, FRR).
      Επαναφέρει τη συχνότητα στην ονομαστική τιμή και τις ανταλλαγές ισχύος μεταξύ περιοχών
      στα προγράμματα και αποδεσμεύει την ΕΔΣ. Διακρίνεται σε δύο μορφές:
      \begin{itemize}
      \item \textbf{αΕΑΣ} (automatic Frequency Restoration Reserve, aFRR), με αυτόματη ενεργοποίηση: κεντρική, οδηγείται από το σφάλμα
            ελέγχου αποκατάστασης συχνότητας (ΣΕΑΣ, ενότητα «\nameref{sec:frr-multi}») μέσω της ΑΡΠ του
            ΔΣΜ·
      \item \textbf{χΕΑΣ} (manual Frequency Restoration Reserve, mFRR), με χειροκίνητη ενεργοποίηση: ενεργοποιείται με εντολή
            κατανομής του ΔΣΜ και αποδεσμεύει την αΕΑΣ σε παρατεταμένες ανισορροπίες.
      \end{itemize}
\item \textbf{Εφεδρεία αντικατάστασης} (ΕΑ· Replacement Reserve, RR). Αποκαθιστά το απαιτούμενο
      επίπεδο των ενεργοποιημένων ΕΑΣ, ώστε το σύστημα να είναι έτοιμο για επόμενη διαταραχή.
\end{itemize}

Στην παλαιότερη βιβλιογραφία και σε μέρος της πρακτικής οι τρεις πρώτες κατηγορίες
αναφέρονται ως πρωτεύουσα, δευτερεύουσα και τριτεύουσα ρύθμιση αντιστοίχως. Η ορολογία του
SO GL, με τα ακρωνύμια της επίσημης ελληνικής απόδοσης που χρησιμοποιεί και ο ΑΔΜΗΕ,
κατονομάζει τις εφεδρείες με βάση τη λειτουργία τους (διατήρηση, αποκατάσταση,
αντικατάσταση) και όχι με βάση τη σειρά ενεργοποίησης.

Κάθε επόμενη κατηγορία εφεδρείας ενεργοποιείται βραδύτερα και με χαμηλότερο κόστος από την
προηγούμενη και την αποδεσμεύει, ώστε οι ταχύτερες εφεδρείες να είναι διαθέσιμες για νέα
διαταραχή. Οι χρόνοι πλήρους ενεργοποίησης δίνονται στον Πίνακα~\ref{tab:reserves}.

\begin{table}[!htbp]
\centering
\begin{tabular}{lllc}
\hline
Εφεδρεία & Λειτουργία & Ενεργοποίηση & Χρόνος πλήρους ενεργοποίησης \\
\hline
ΕΔΣ (FCR)   & διατήρηση      & αυτόματη, τοπική   & 30 s \\
αΕΑΣ (aFRR) & αποκατάσταση   & αυτόματη, κεντρική & 5 min \\
χΕΑΣ (mFRR) & αποκατάσταση   & χειροκίνητη        & 12,5 min \\
ΕΑ (RR)     & αντικατάσταση  & χειροκίνητη        & $>$ 15 min \\
\hline
\end{tabular}
\caption{Επίπεδα εφεδρείας συχνότητας στην Ηπειρωτική Ευρώπη, κατά τον SO GL και τα
τυποποιημένα προϊόντα εξισορρόπησης.}
\label{tab:reserves}
\end{table}

Οι επόμενες ενότητες εξετάζουν διαδοχικά την αδράνεια, την ΕΔΣ και την ΕΑΣ σε απομονωμένο
και σε διασυνδεδεμένο σύστημα. Το Σχήμα~\ref{fig:freq} συνοψίζει την ακολουθία για το συμβάν
αναφοράς της Ηπειρωτικής Ευρώπης.

\section{Αδράνεια και αρχικός ρυθμός μεταβολής συχνότητας}

Σε μόνιμη κατάσταση η συχνότητα είναι σταθερή και κοινή σε ολόκληρη τη συγχρονισμένη
περιοχή· κάθε ανισορροπία μεταξύ παραγωγής και ζήτησης προκαλεί μεταβολή της.

Αμέσως μετά από ανισορροπία ισχύος $\Delta P$, για παράδειγμα την απώλεια μεγάλης μονάδας
ηλεκτροπαραγωγής, οι ρυθμιστές στροφών δεν έχουν ακόμη αποκριθεί και το έλλειμμα καλύπτεται
από την κινητική ενέργεια των στρεφόμενων μαζών των σύγχρονων μηχανών. Ο αρχικός ρυθμός
μεταβολής της συχνότητας (Rate of Change of Frequency, RoCoF) προκύπτει από την εξίσωση
ταλάντωσης (swing equation) και δίνεται κατ' απόλυτη τιμή από τη σχέση

\begin{equation}
\frac{\mathrm{d}f}{\mathrm{d}t} = \frac{f_0 \, \Delta P}{2 H S_n}
\label{eq:rocof}
\end{equation}

όπου $f_0$ η ονομαστική συχνότητα, $H$ η σταθερά αδράνειας του συστήματος σε δευτερόλεπτα
και $S_n$ η συνολική ονομαστική φαινόμενη ισχύς των \emph{συγχρονισμένων} μονάδων.

\emph{Παράδειγμα.} Για συγχρονισμένη περιοχή με $S_n = 300$ GVA, $H = 5$ s και απώλεια
$\Delta P = 3\,000$ MW, ο αρχικός ρυθμός μεταβολής είναι
$50 \times 3000 / (2 \times 5 \times 300\,000)$ = 0,05 Hz/s. Για την ίδια απώλεια σε
συνθήκες χαμηλού φορτίου, με το ήμισυ της συγχρονισμένης ισχύος, ο ρυθμός διπλασιάζεται.

Σύμφωνα με την \eqref{eq:rocof}, η αντικατάσταση σύγχρονων γεννητριών από μονάδες που
συνδέονται μέσω ηλεκτρονικών μετατροπέων ισχύος, όπως τα αιολικά και φωτοβολταϊκά πάρκα,
μειώνει την αποθηκευμένη κινητική ενέργεια $H S_n$ και αυξάνει τον ρυθμό μεταβολής της
συχνότητας. Μειώνεται έτσι ο διαθέσιμος χρόνος πριν η συχνότητα φθάσει στο κατώφλι της
αυτόματης αποσύνδεσης ζήτησης λόγω χαμηλής συχνότητας, και απαιτούνται ταχύτερες υπηρεσίες,
όπως η συνθετική αδράνεια και η ταχεία απόκριση συχνότητας.

Η μείωση της αδράνειας είναι ήδη μετρήσιμη. Για την Ηπειρωτική Ευρώπη, ο ENTSO-E εκτίμησε την
αδράνεια του 2019 ανά ώρα από το μείγμα παραγωγής και παρουσίασε την κατανομή της ανά μήνα και
ώρα της ημέρας, με εμφανή μείωση τις μεσημεριανές ώρες λόγω της φωτοβολταϊκής παραγωγής, η
οποία αναμένεται να ενταθεί έως το 2030 (Σχήμα~\ref{fig:inertia-ce}). Το 2019 η σταθερά
αδράνειας του συστήματος κυμαινόταν κυρίως μεταξύ 4,5 και 6 s, με τις χαμηλότερες τιμές τις
μεσημεριανές ώρες, ενώ στο σενάριο του 2030 οι μεσημεριανές τιμές της άνοιξης και του
καλοκαιριού μειώνονται περίπου στα 1,5 έως 3 s. Ο SO GL προβλέπει κοινή μελέτη των ΔΣΜ κάθε
συγχρονισμένης περιοχής, η οποία επικαιροποιείται κάθε δύο έτη, για να διαπιστωθεί αν πρέπει
να καθοριστούν απαιτήσεις ελάχιστης αδράνειας.

\begin{figure}[!htbp]
\centering
\includegraphics[width=\linewidth]{figures/entsoe-inertia-2019-2030.png}
\caption{Κατανομή της σταθεράς αδράνειας $H_{\mathrm{sys}}$ της Ηπειρωτικής Ευρώπης ανά μήνα
(οριζόντιος άξονας) και ώρα της ημέρας (κατακόρυφος άξονας), εκτιμώμενη από το μείγμα
παραγωγής, για το 2019 (αριστερά) και για το σενάριο Distributed Generation 2030 του δεκαετούς προγράμματος ανάπτυξης δικτύου (Ten-Year
Network Development Plan, TYNDP) 2018
(δεξιά). Πηγή: ENTSO-E,
\href{https://eepublicdownloads.entsoe.eu/clean-documents/SOC\%20documents/Inertia\%20and\%20RoCoF_v17_clean.pdf}{«Inertia
and Rate of Change of Frequency (RoCoF)»}, έκδοση 17, 2020, Σχήμα 10.}
\label{fig:inertia-ce}
\end{figure}

Το Σχήμα~\ref{fig:inertia} δείχνει την εξέλιξη σε ένα μικρότερο σύστημα, της Μεγάλης Βρετανίας,
για το οποίο ο διαχειριστής NESO (National Energy System Operator) δημοσιεύει ως ανοιχτά δεδομένα την αδράνεια ανά μισάωρο.
Μεταξύ 2019 και 2025 η μέση αδράνεια μειώθηκε από 241 σε 178 GVA·s, δηλαδή κατά 26\% περίπου,
το ποσοστό του χρόνου με αδράνεια κάτω από 150 GVA·s αυξήθηκε από 2\% σε 37\% και η ελάχιστη
τιμή μειώθηκε από 135 σε 108 GVA·s. Οι χαμηλότερες τιμές εμφανίζονται τις νυχτερινές ώρες και,
το 2025, και τις μεσημεριανές ώρες από την άνοιξη έως το φθινόπωρο, όταν μεγάλο μέρος της
ζήτησης καλύπτεται από αιολική και φωτοβολταϊκή παραγωγή. Με βάση την \eqref{eq:rocof}, για
το ίδιο συμβάν ο αρχικός ρυθμός μεταβολής της συχνότητας είναι αντιστρόφως ανάλογος της
αδράνειας: στην ελάχιστη τιμή του 2025 είναι περίπου 2,2 φορές μεγαλύτερος από ό,τι στη μέση
τιμή του 2019.

\begin{figure}[!htbp]
\centering
\includegraphics[width=\linewidth]{figures/inertia-heatmap-gb.svg}
\caption{Αδράνεια του συστήματος της Μεγάλης Βρετανίας ανά μισάωρο (τοπική ώρα) για τα έτη 2019
και 2025. Κάθε κατακόρυφη λωρίδα αντιστοιχεί σε μία ημέρα και οι ανοιχτότεροι χρωματισμοί σε
χαμηλότερη αδράνεια. Δεδομένα:
\href{https://www.neso.energy/data-portal/system-inertia}{NESO, System Inertia}, Outturn Inertia.
Supported by National Energy SO Open Data.}
\label{fig:inertia}
\end{figure}

Το Σχήμα~\ref{fig:freq} απεικονίζει την απόκριση συχνότητας για το συμβάν αναφοράς και τις
τρεις διαδοχικές φάσεις της.

\begin{figure}[!htbp]
\centering
\includegraphics[width=\linewidth]{figures/frequency-response.svg}
\caption{Απόκριση συχνότητας σε απώλεια παραγωγής 3\,000 MW. Η αρχική κλίση καθορίζεται
αποκλειστικά από την αδράνεια· η εφεδρεία διατήρησης συχνότητας (ΕΔΣ) ανακόπτει την πτώση
και σταθεροποιεί τη συχνότητα σε νέα τιμή μόνιμης κατάστασης· η εφεδρεία αποκατάστασης
συχνότητας (ΕΑΣ) εξαλείφει το μόνιμο σφάλμα και επαναφέρει τη συχνότητα στα 50 Hz. Η
διακεκομμένη καμπύλη είναι η απόκριση μόνο με ΕΔΣ.}
\label{fig:freq}
\end{figure}

\section{Εφεδρεία διατήρησης συχνότητας και στατισμός}

Η ΕΔΣ υλοποιείται από τους ρυθμιστές στροφών (governors) των μονάδων, οι οποίοι μεταβάλλουν
την παραγόμενη ισχύ αναλογικά προς την απόκλιση συχνότητας. Η αναλογία καθορίζεται από τον
\textbf{στατισμό} (droop) $R$ της μονάδας, ο οποίος ορίζεται ως ο λόγος της ανά μονάδα
μεταβολής της συχνότητας προς την ανά μονάδα μεταβολή της ισχύος εξόδου σε μόνιμη κατάσταση,
με βάσεις την ονομαστική συχνότητα $f_0$ και την ονομαστική ισχύ $P_n$ της μονάδας:

\begin{equation}
R = -\frac{\Delta f / f_0}{\Delta P / P_n}
\label{eq:droopdef}
\end{equation}

ή ισοδύναμα

\begin{equation}
\frac{\Delta P}{P_n} = -\frac{1}{R}\,\frac{\Delta f}{f_0}
\label{eq:droop}
\end{equation}

Ο στατισμός είναι η κλίση της στατικής χαρακτηριστικής συχνότητας-ισχύος της μονάδας και
εκφράζεται συνήθως επί τοις εκατό. Στατισμός 5\% σημαίνει ότι μεταβολή της συχνότητας κατά 5\%
της ονομαστικής, δηλαδή κατά 2,5 Hz σε σύστημα 50 Hz, μεταβάλλει την ισχύ εξόδου κατά την
ονομαστική ισχύ της μονάδας· μικρότερος στατισμός σημαίνει ισχυρότερη απόκριση. Το αντίστροφο
$1/R$ είναι το κέρδος του αναλογικού ρυθμιστή. Με μηδενικό στατισμό (ισόχρονη ρύθμιση) η
μονάδα θα επανέφερε μόνη της τη συχνότητα στην ονομαστική τιμή, αλλά περισσότερες ισόχρονες
μονάδες δεν μπορούν να λειτουργούν παράλληλα, επειδή η κατανομή του φορτίου μεταξύ τους δεν
είναι καθορισμένη.

Σε μόνιμη κατάσταση όλες οι μονάδες της συγχρονισμένης περιοχής αντιλαμβάνονται την ίδια
απόκλιση συχνότητας, οπότε από την \eqref{eq:droop} η συνεισφορά της μονάδας $i$ είναι
$\Delta P_i = -P_{n,i}\,\Delta f / (R_i f_0)$: μονάδες με ίσο στατισμό μοιράζονται τη
μεταβολή φορτίου αναλογικά προς την ονομαστική τους ισχύ. Ο στατισμός καθορίζει επομένως τη
συμμετοχή κάθε μονάδας χωρίς ανάγκη επικοινωνίας μεταξύ τους. Στην Ηπειρωτική Ευρώπη η ΕΔΣ
ενεργοποιείται πλήρως για απόκλιση συχνότητας $\pm 200$ mHz, εντός 30 s, ενώ η συνολική νεκρή
ζώνη και αναισθησία της απόκρισης δεν επιτρέπεται να υπερβαίνει τα 10 mHz.

Τυπικές τιμές στατισμού είναι $R = 4$ έως 5\%. Για παράδειγμα, μονάδα 500 MW με $R = 5\%$,
για απόκλιση $\Delta f = -200$ mHz, δηλαδή −0,4\%, αυξάνει την παραγωγή της κατά
(0,004~/~0,05)~×~500~=~40~MW.

Το αρνητικό πρόσημο αντιστοιχεί σε αρνητική ανάδραση: μείωση της συχνότητας προκαλεί αύξηση
της παραγωγής. Επειδή ο έλεγχος είναι \emph{αναλογικός}, η πρόσθετη παραγωγή διατηρείται μόνο
όσο διατηρείται η απόκλιση συχνότητας. Για ανισορροπία $\Delta P$ η απόκλιση μόνιμης
κατάστασης είναι $\Delta f_{\mathrm{ss}} = -\Delta P / \lambda$, όπου $\lambda$ (MW/Hz) η
χαρακτηριστική ισχύος-συχνότητας της περιοχής, δηλαδή το άθροισμα της συνεισφοράς των
ρυθμιστών και της αυτορρύθμισης του φορτίου. Η ΕΔΣ επομένως σταθεροποιεί τη συχνότητα σε
νέα τιμή μόνιμης κατάστασης, με μόνιμο σφάλμα ως προς την ονομαστική τιμή, το οποίο
εξαλείφεται από την εφεδρεία αποκατάστασης συχνότητας.

\section{Εφεδρεία αποκατάστασης συχνότητας σε απομονωμένο σύστημα}

Η ΕΑΣ εξετάζεται αρχικά για \emph{απομονωμένο} σύστημα, χωρίς διασυνδέσεις με γειτονικά
συστήματα. Στην περίπτωση αυτή το μόνο μέγεθος ελέγχου είναι η απόκλιση συχνότητας και ο
κεντρικός ελεγκτής εφαρμόζει \emph{ολοκληρωτική} δράση:

\begin{equation}
\Delta P_{\mathrm{FRR}}(t) = -K_I \int_0^{t} \Delta f(\tau)\,\mathrm{d}\tau
\label{eq:frr-isolated}
\end{equation}

Όσο $\Delta f \neq 0$, το ολοκλήρωμα μεταβάλλεται και η εντολή παραγωγής αυξάνεται· σε μόνιμη
κατάσταση ισχύει κατ' ανάγκη $\Delta f = 0$. Ο ολοκληρωτικός έλεγχος εξαλείφει επομένως το
μόνιμο σφάλμα της αναλογικής ΕΔΣ. Με την επαναφορά της συχνότητας στην ονομαστική τιμή, η
συνεισφορά της ΕΔΣ μηδενίζεται σύμφωνα με την \eqref{eq:droop}, δηλαδή η ΕΑΣ
\emph{αποδεσμεύει} την ΕΔΣ, η οποία καθίσταται εκ νέου διαθέσιμη για επόμενη διαταραχή.

Η απόκριση της ΕΑΣ σχεδιάζεται σημαντικά βραδύτερη από εκείνη της ΕΔΣ (λεπτά έναντι
δευτερολέπτων). Ο διαχωρισμός των χρονικών κλιμάκων αποτρέπει την αλληλεπίδραση των δύο
βρόχων ελέγχου και τις ταλαντώσεις που θα προέκυπταν από αυτήν.

\section{Εφεδρεία αποκατάστασης συχνότητας σε σύστημα πολλών περιοχών}
\label{sec:frr-multi}

Σε διασυνδεδεμένο σύστημα η συχνότητα είναι κοινή σε ολόκληρη τη συγχρονισμένη περιοχή,
οπότε απώλεια παραγωγής σε οποιαδήποτε περιοχή προκαλεί απόκλιση συχνότητας σε όλες. Αν ο
ελεγκτής κάθε περιοχής χρησιμοποιούσε μόνο την απόκλιση συχνότητας, όλες οι περιοχές θα
ενεργοποιούσαν ΕΑΣ για την ίδια διαταραχή. Αυτό θα προκαλούσε υπερδιόρθωση, αποκλίσεις των
ανταλλαγών ισχύος από τα προγράμματα και μετακύλιση του κόστους εξισορρόπησης στις περιοχές
που δεν προκάλεσαν τη διαταραχή.

Για τον διαχωρισμό της περιοχής στην οποία εμφανίστηκε η ανισορροπία, ο ελεγκτής κάθε
περιοχής ελέγχου φορτίου-συχνότητας (περιοχή ΕΦΣ· Load-Frequency Control area, LFC area) μηδενίζει το \emph{σφάλμα ελέγχου
περιοχής} (ΣΕΠ· Area Control Error, ACE):

\begin{equation}
\mathrm{ACE} = \Delta P_{\mathrm{tie}} + K \, \Delta f
\label{eq:ace}
\end{equation}

όπου $\Delta P_{\mathrm{tie}}$ το σφάλμα ελέγχου ισχύος, δηλαδή η απόκλιση της μετρούμενης
ανταλλαγής ισχύος στις γραμμές διασύνδεσης από το πρόγραμμα, και $K$ ο συντελεστής K της
περιοχής σε MW/Hz (στην κλασική βιβλιογραφία συντελεστής πόλωσης συχνότητας, frequency
bias). Στο απομονωμένο σύστημα ο πρώτος όρος δεν υπάρχει και η \eqref{eq:ace} ανάγεται στο
κριτήριο της προηγούμενης ενότητας. Για τον λόγο αυτόν ο SO GL ορίζει ως ενιαίο μέγεθος
ελέγχου της αποκατάστασης το \emph{σφάλμα ελέγχου αποκατάστασης συχνότητας} (ΣΕΑΣ·
Frequency Restoration Control Error, FRCE): σε σύστημα πολλών περιοχών ισούται με το ΣΕΠ
της περιοχής, ενώ όταν η περιοχή ελέγχου συμπίπτει με ολόκληρη τη συγχρονισμένη περιοχή
ισούται με την απόκλιση συχνότητας.

Ο συνδυασμός των δύο όρων διακρίνει τη θέση της διαταραχής (Σχήμα~\ref{fig:lfc}):

\begin{itemize}
\item \textbf{Διαταραχή εντός της περιοχής.} Η συχνότητα μειώνεται ($\Delta f < 0$) και οι
      γειτονικές περιοχές παρέχουν ισχύ υποστήριξης, οπότε οι εισαγωγές υπερβαίνουν το
      πρόγραμμα ($\Delta P_{\mathrm{tie}} < 0$). Οι δύο όροι είναι ομόσημοι, το ΣΕΠ λαμβάνει
      μεγάλη απόλυτη τιμή και ενεργοποιείται ΕΑΣ στην περιοχή.
\item \textbf{Διαταραχή εκτός της περιοχής.} Η συχνότητα μειώνεται, αλλά η περιοχή εξάγει
      πρόσθετη ισχύ προς τη γειτονική περιοχή ($\Delta P_{\mathrm{tie}} > 0$). Για κατάλληλη
      τιμή του $K$ οι δύο όροι αλληλοαναιρούνται και το ΣΕΠ παραμένει σχεδόν μηδενικό· η
      περιοχή συμμετέχει μόνο μέσω της ΕΔΣ.
\end{itemize}

Η ιδιότητα αυτή ονομάζεται \emph{μη αλληλεπιδραστικός έλεγχος} (non-interactive control):
κάθε περιοχή αντισταθμίζει μόνο τις δικές της ανισορροπίες, ενώ όλες οι περιοχές
συμμετέχουν μέσω της ΕΔΣ στη συγκράτηση της συχνότητας αμέσως μετά τη διαταραχή. Η
αλληλοαναίρεση είναι πλήρης όταν ο συντελεστής $K$ ισούται με τη χαρακτηριστική
ισχύος-συχνότητας της περιοχής, δηλαδή με το άθροισμα της συνεισφοράς της ΕΔΣ, του
αυτοελέγχου της παραγωγής και της αυτορρύθμισης του φορτίου ανά Hz.

\begin{figure}[!htbp]
\centering
\includegraphics[width=\linewidth]{figures/lfc-areas.svg}
\caption{Εφεδρεία αποκατάστασης συχνότητας σε απομονωμένο και σε διασυνδεδεμένο σύστημα. Στην
πρώτη περίπτωση αρκεί η απόκλιση συχνότητας· στη δεύτερη απαιτείται και ο όρος των
ανταλλαγών ισχύος στις γραμμές διασύνδεσης, ώστε να ενεργοποιείται ΕΑΣ μόνο στην περιοχή
όπου εμφανίστηκε η διαταραχή.}
\label{fig:lfc}
\end{figure}

\section{Παράμετροι συχνότητας για την Ηπειρωτική Ευρώπη}

\begin{itemize}
\item Το \textbf{συμβάν αναφοράς} (reference incident) για τη διαστασιολόγηση της ΕΔΣ είναι
      ανισορροπία 3\,000 MW, είτε ως απώλεια παραγωγής είτε ως απώλεια φορτίου. Η ελάχιστη
      διαθέσιμη ΕΔΣ στη συγχρονισμένη περιοχή είναι ανά πάσα στιγμή 3\,000 MW.
\item Η \textbf{μέγιστη απόκλιση συχνότητας σταθερής κατάστασης} για το συμβάν αναφοράς
      είναι 200 mHz. Από τα δύο μεγέθη προκύπτει ο συντελεστής αναλογίας μεταξύ ΕΔΣ και
      απόκλισης συχνότητας: 3\,000 MW / 0,2 Hz = 15\,000 MW/Hz.
\item Η \textbf{τυπική περιοχή συχνότητας} (standard frequency range) είναι $\pm 50$ mHz· η
      συχνότητα δεν επιτρέπεται να βρίσκεται εκτός αυτής για περισσότερα από
      15\,000 λεπτά ανά έτος. Η μέγιστη στιγμιαία απόκλιση συχνότητας είναι 800 mHz.
\item Απόκλιση συχνότητας σταθερής κατάστασης μεγαλύτερη των 200 mHz κατατάσσει το σύστημα
      σε κατάσταση \emph{έκτακτης ανάγκης} (Κεφάλαιο 1, ενότητα
      «\nameref{sec:states}»).
\end{itemize}

% TODO: παραδείγματα και λυμένες ασκήσεις (σελίδα 03b-paradeigmata.md, όπως στο Κεφάλαιο 1)
